\chapter{Literature Review}

\section{Comprehensive Literature Review}

\subsection{Evolution of Governance, Risk, and Compliance (GRC) Platforms}
Historical security auditing relied on point-in-time assessments, leaving temporal blind spots between evaluation periods. Organizations managed Governance, Risk, and Compliance (GRC) through manual evidence collection, rendering the compliance state obsolete following an assessment. Regulatory frameworks require a transition toward continuous controls monitoring architectures. Standards including NIST Special Publication 800-53 Revision 5 and ISO/IEC 27001:2022 mandate ongoing monitoring to validate control implementation \cite{nist80053_2020, iso27001_2022}. The Federal Information Security Modernization Act (FISMA) and the Federal Risk and Authorization Management Program (FedRAMP) require agencies to automate the tracking of systemic anomalies \cite{fisma_2014}. The transition to SIEM-driven compliance introduces operational advantages. Centralized log management supports audit accountability by preserving data confidentiality, integrity, and availability \cite{nist80092_2006}. Automated analysis establishes operational baselines, facilitating forensic investigations during security incidents. Security administrators configure SIEM agents to track software inventories, hardware configurations, and system health telemetry, feeding this data into the GRC database. Continuous validation replaces reactive audits with a data-driven security posture. By evaluating log streams against established baselines, organizations ensure evidence generation becomes a byproduct of daily security operations.

\subsection{Security Information and Event Management (SIEM) Integrations}
The integration of Security Information and Event Management (SIEM) systems alters the GRC methodology. SIEM platforms act as telemetry ingestion engines, replacing manual log retrieval with automated data correlation. NIST SP 800-53 defines control families, notably Audit and Accountability (AU) and System and Information Integrity (SI), which require automated tracking to satisfy continuous validation requirements \cite{nist80053_2020}. Automating these controls involves mapping endpoint telemetry directly to regulatory requirements. Platforms like Wazuh, an open-source SIEM and Extended Detection and Response (XDR) solution, facilitate this translation by supplying event aggregation logic that maps compliance events against NIST and ISO matrices \cite{wazuh_doc}. Automated compliance mapping bridges the semantic gap between technical events and business risk. When a user modifies a protected configuration file, File Integrity Monitoring (FIM) sensors generate an event log. The SIEM correlates this event against framework matrices, flagging the action as a violation of ISO 27001 Annex A.8.9 (Configuration Management) or NIST SP 800-53 control CM-6 \cite{nist800128_2011}. Analysts monitor deviations via dashboards that aggregate alert volumes, isolating non-compliant agents for remediation. The Wazuh Security Configuration Assessment (SCA) module provides the deterministic validation required to satisfy these governance frameworks by evaluating configuration states against security benchmarks.

\subsection{Performance Bottlenecks in Telemetry Aggregation}
Calculating compliance metrics across distributed fleets exposes data aggregation limitations in traditional GRC platforms. Standard relational database implementations suffer from the N+1 query problem, where the application executes sequential retrieval loops for individual security controls. This architectural flaw imposes a linear $O(N)$ algorithmic complexity on data retrieval, causing geometric performance decay as monitored agents increase \cite{planetscale_n1}. The resulting latency prevents organizations from achieving continuous monitoring, because dashboard rendering times exceed thresholds for operational visibility. The overhead associated with Object-Relational Mapping (ORM) frameworks extends beyond query volume. Application-level data processing, where the system retrieves full object graphs before filtering or aggregating in the Python memory space, creates memory bandwidth bottlenecks. This pattern forces the application server to instantiate transient objects, triggering garbage collection cycles that block the execution of compliance logic. In environments where SIEM agents transmit telemetry continuously, the cumulative impact of these $O(N)$ operations leads to a loss of operational visibility.

Optimizing relational data access requires transitioning from sequential $O(N)$ looping to $O(1)$ set-based aggregations. Set-based modeling formalizes relational dynamics, demonstrating that task-informed database representations retain relevant signals more efficiently than application-level iteration \cite{kuraj_microservices_2024}. Developers optimize ORM queries by utilizing eager loading mechanisms that translate application logic into SQL JOIN clauses. Fetching related entities in a single query allows the database engine's internal query optimizer to utilize B-Tree indices efficiently, reducing total execution time and minimizing context switching between the application runtime and the database server. Set-based database aggregations shift the mathematical workload to the relational engine. Using query projection techniques, the system instructs the database to compute sums, averages, and conditional compliance counts internally. This methodology reduces data transfer across the network wire from megabytes of row data to primitive integer values. By utilizing conditional aggregation functions, where the engine evaluates compliance criteria during the initial table scan, the platform achieves a deterministic $O(1)$ response time. Implementing these optimizations ensures dashboard rendering remains responsive as the organizational asset inventory scales. The architectural decision to prioritize database-side computation over application-side iteration reflects the principle of pushing logic closer to the data. This approach ensures the GRC platform can process telemetry logs without exhausting the system's computational resources.

\subsection{Modern Access Control Paradigms}
Securing the compliance data layer against unauthorized extraction requires Identity and Access Management (IAM) protocols. Role-Based Access Control (RBAC) provides the foundation for authorization within distributed systems. RBAC methodologies control access to system resources based on roles assigned to a user, rather than assigning explicit permissions directly to the user identity \cite{nelson_stateless_2024}. Authentication in distributed microservices relies on stateless session security, primarily implemented via JSON Web Tokens (JWT). The authorization server cryptographically signs the JWT using asymmetric encryption algorithms, preventing payload tampering and allowing local verification without initiating a database lookup. Authorization in the ICMS platform relies on an enforcement layer that integrates Zero Trust principles into the API routing logic. Rather than assuming internal network trust, the system evaluates the identity context of every inbound request. Custom permission classes intercept traffic at the view layer, executing Boolean logic to verify the requester possesses the RBAC attributes required for the target resource. This enforcement remains decoupled from the business logic, ensuring a consistent security boundary across the microservice ecosystem.

The security lifecycle of stateless JSON Web Tokens (JWT) incorporates rotation policies to mitigate risks inherent in long-lived credentials. Traditional session IDs require database lookups to verify validity; JWTs carry cryptographically signed claims that allow microservices to authenticate users locally. To prevent session hijacking via intercepted tokens, the platform implements a refresh token rotation scheme. When the client requests a new access credential, the server invalidates the previous refresh token and issues a new pair, ensuring compromised credentials have limited utility. The architecture incorporates Time-Based One-Time Passwords (TOTP) as a secondary authentication factor. TOTP algorithms rely on a high-entropy, shared cryptographic secret established between the server and the client's authenticator application. Both entities use the current Unix timestamp, divided by a predetermined time interval, as the moving factor in a Keyed-Hash Message Authentication Code (HMAC-SHA1) function \cite{rfc6238}. Because the resulting 6-digit numeric code mutates deterministically every 30 seconds, intercepted credentials become mathematically useless. Secondary authentication factors establish an additional defensive layer that remains resilient against replay attacks. Client-side interceptors monitor for specific HTTP response states; when the server returns a forbidden status due to a missing or invalid TOTP challenge, the application redirects the user to a step-up authentication view, preventing unauthorized access. This session lifecycle ensures the compliance data remains inaccessible to compromised endpoints.

\subsection{Containerization and Defense-in-Depth Topology}
Enterprise architecture increasingly transitions from monolithic software design toward microservices and containerized deployments. Monolithic systems couple application logic, user interfaces, and database drivers into a single deployment artifact, where a failure in one component frequently cascades across the runtime environment. Containerization technologies, specifically Docker, provide the standardized runtime environment for these distributed microservices, ensuring deterministic scalability and bounded contexts through Linux kernel namespaces and control groups. Deploying the ICMS components via isolated Docker virtual bridge networks establishes a defense-in-depth topology. The architecture physically separates the public-facing API gateway layers from the private persistent storage layers, such as PostgreSQL and Redis. By binding these data tiers to internal, non-routable Docker bridge networks, the system remains inaccessible to external network scans. This isolation ensures actors cannot establish direct TCP connections with the database, forcing data interactions through the authenticated business-logic layer. Adherence to external hardening standards, such as NIST SP 800-190 and the CIS Docker Benchmark, further neutralizes execution environment vulnerabilities \cite{nist800190_2017, cis_docker_2020}.

\subsection{Cognitive Load and HCI in Security Dashboards}
Security Operations Centers process large volumes of diagnostic data daily, frequently triggering alert fatigue. Legacy GRC platforms present analysts with unfiltered data in dense HTML tables, leading to sensory overload and an increase in diagnostic errors. High cognitive load correlates with delayed incident response times and operator burnout \cite{mit_hci_security_2022}. Mitigating cognitive strain requires the application of User Interface (UI) design principles based on human-computer interaction studies. Modern browser rendering engines deprecate structural tables for layout purposes in favor of CSS Flexbox and Grid modules, which allow child components to wrap and scale fluidly based on viewport dimensions \cite{mdn_flexbox}. The ICMS platform utilizes modular card components and semantic color scaling, such as bold red for critical exploits and minimal grayscale for benign logs, to guide the analyst's focus toward critical threats. Leveraging elevation through depth illusions distinguishes interactive components from static background data, structuring telemetry into a decision-support matrix.

\section{Analysis of Existing Systems}

\subsection{Manual Spreadsheet Auditing}
Evaluating the baseline of traditional governance reveals operational limitations inherent in manual compliance methodologies. Historically, organizations relied on labor-intensive processes encompassing decentralized spreadsheets, task-tracking tickets, and point-in-time audits conducted by external consultants. This fragmented approach fails to scale within IT environments, because the compliance state captured during an audit becomes obsolete upon the conclusion of the assessment. The reliance on manual sampling and evidence collection introduces temporal blind spots, where vulnerabilities remain undetected between audit cycles. The structural limits of spreadsheet applications lead to performance degradation and data corruption when tracking large volumes of configuration parameters. Manual mapping processes introduce a high probability of human error, as analysts attempt to interpret technical data against regulatory constraints.

\subsection{Vanta (Commercial SaaS)}
Vanta operates in the commercial Software-as-a-Service (SaaS) continuous monitoring sector. The platform accelerates the evidence collection process by deploying proprietary agents and utilizing API connectors to integrate with public cloud providers and enterprise identity systems. This commercial solution introduces risks regarding data sovereignty and vendor lock-in. Operating as a multi-tenant cloud environment, Vanta requires client organizations to transmit security telemetry to its external, third-party infrastructure. This data exfiltration increases the attack surface, introducing potential exposure via shared infrastructure vulnerabilities. Vanta operates as a proprietary system; its closed-source algorithms prevent client organizations from independently verifying or customizing how raw telemetry events map to compliance controls. This lack of transparency forces organizations to rely on the provider's internal security assertions, contradicting the principles of algorithmic transparency.

\subsection{Drata (Commercial SaaS)}
Drata provides automated compliance tracking via a centralized, multi-tenant SaaS architecture. While the platform offers integrations and pre-mapped framework tracking, it involves risks regarding external data storage. Organizations utilizing Drata export their infrastructure metadata to the provider's cloud, violating data sovereignty requirements mandated by entities in the defense and financial sectors. The closed-source nature of Drata restricts security engineers from inspecting the underlying evaluation logic. This limitation prevents custom programmatic modifications to the compliance engine, meaning that organizations with unique internal architectures or non-standard security tools cannot extend the system without relying on proprietary vendor interventions. The inability to inspect the evaluation engine ensures the platform remains a closed-box compliance oracle rather than a verifiable security tool.

\subsection{AuditBoard (Enterprise GRC)}
AuditBoard targets the enterprise GRC market, focusing on risk management, ESG reporting, and internal audit synchronization. Despite its feature set, AuditBoard introduces integration rigidity. The platform demands high licensing costs and requires professional service engagements for initial deployment, resulting in deployment cycles that frequently span multiple financial quarters. This enterprise architecture compares poorly against containerized solutions that deploy rapidly. AuditBoard relies on proprietary connectors to interface with security tools, establishing a closed-box mapping approach. Security engineers lack the capacity to manipulate the programmatic logic linking SIEM alerts to organizational risk matrices. This opacity limits the platform's utility in environments that demand deterministic, verifiable compliance states.

\section{Comparative Analysis}

The ICMS platform establishes an alternative for organizations requiring data sovereignty, technical transparency, and verifiable execution logic. By integrating with the open-source Wazuh SIEM, the platform achieves automation without necessitating external data transmission. The architecture provides programmatic control via the open-source Django framework, allowing organizations to independently audit, modify, and extend the underlying compliance evaluation algorithms. The single-tenant Docker topology ensures that security telemetry remains within an isolated, customer-controlled bridge network, neutralizing the attack surface associated with multi-tenant cloud environments. 

As detailed in Table \ref{tab:comparative_matrix}, the ICMS addresses the performance decay observed in proprietary systems through database engineering. By implementing $O(1)$ set-based aggregations via query annotation and SQL filtering, the platform processes telemetry logs efficiently. This deterministic execution ensures the platform maintains lower latency compared to legacy systems burdened by monolithic database scaling constraints.

    \begin{table}[htbp]
        \centering
        \caption{Comparative Feature Matrix: Analyzing GRC Platforms}
        \label{tab:comparative_matrix}
        \renewcommand{\arraystretch}{1.5}
        \small
        \begin{tabular}{|>{\raggedright\arraybackslash}p{2.8cm}|>{\raggedright\arraybackslash}p{3.5cm}|>{\raggedright\arraybackslash}p{4.3cm}|>{\raggedright\arraybackslash}p{3.5cm}|}
            \hline
            \textbf{Feature} & \textbf{Traditional Methods (Manual)} & \textbf{Commercial SaaS (Vanta / Drata / AuditBoard)} & \textbf{ICMS Platform} \\
            \hline
            Deployment Model & Decentralized documents & Multi-tenant public cloud (SaaS risk) & Single-tenant, isolated Docker bridge \\
            \hline
            Data Sovereignty & High (Data remains local) & Low (Data transmitted to third parties) & High (Data confined to internal network) \\
            \hline
            Performance Complexity & Prone to spreadsheet limits and crashes & Proprietary ORM scaling constraints, API latency & $O(1)$ set-based aggregation via .annotate() \\
            \hline
            Customizability & Manual modification & Rigid black-box mapping algorithms & Full programmatic control via Django ORM \\
            \hline
            Cost & High operational labor expenditure & Expensive recurring enterprise licensing & Zero licensing cost (Open-source stack) \\
            \hline
        \end{tabular}
    \end{table}

\section{Summary of Competitive Advantages}

The ICMS platform provides operational and technical advantages that address the flaws of legacy auditing and commercial SaaS offerings:

\begin{itemize}
    \item \textbf{Data Sovereignty via Docker Bridge Isolation:} The platform confines all PostgreSQL database interactions to a private virtual network, ensuring compliance telemetry remains isolated from the public internet and external third-party vendors.
    \item \textbf{Transparent, Verifiable Mapping:} By utilizing an open-source Django backend, organizations can programmatically inspect and verify the deterministic logic linking raw Wazuh SIEM alerts to governance controls, eliminating the proprietary closed-box approach.
    \item \textbf{Deterministic $O(1)$ Database Performance:} Set-based aggregations replace standard ORM looping structures, allowing the system to render fleet-wide compliance dashboards with low latency regardless of the agent inventory size.
    \item \textbf{Zero-Cost Licensing Architecture:} The operational stack relies on open-source technologies (Wazuh, PostgreSQL, Python), removing the financial burden of recurring commercial SaaS fees and professional deployment services.
    \item \textbf{Automated Telemetry Deduplication:} The implementation of a Worst-Case Priority algorithm ensures that conflicting telemetry reports from disparate sensors are resolved deterministically, prioritizing organizational safety and preventing false-positive compliance states.
\end{itemize}
